% Origen: Secciones/PI - Verificación de normativa nacional e internacional.md
\chapter{Verificación de normativa nacional e internacional}
\label{C:normativa}

% Figuras de ESTE capítulo + las compartidas. Siempre al principio del archivo.
\graphicspath{{figs/cap4/}{figs/comunes/}}

En este capítulo se analiza qué normativa nacional e internacional alcanza al
MCA, en qué carácter (obligación, guía de diseño o buena práctica) y qué parte
del trabajo responde a cada una.

\section{Análisis de ARN}
\label{S:norm-arn}

En Argentina, el único marco regulatorio nacional para la actividad nuclear es
el de la Autoridad Regulatoria Nuclear (ARN). La Ley N\textsuperscript{o} 24.804
le asigna la regulación de la seguridad radiológica y nuclear, y la faculta a
dictar normas de cumplimiento obligatorio~\cite{Ley24804}. Sus normas (serie AR)
son públicas. Para este trabajo son relevantes dos: la de diseño de reactores de
investigación, AR~4.2.2~\cite{ARN_AR422}, y la de sistemas de gestión,
AR~10.6.1~\cite{ARN_AR1061}. La ARN no tiene una norma específica para
instrumentación y control digital, software ni FPGA. Su enfoque está basado en
el desempeño: el responsable demuestra que cumple los objetivos de seguridad y
elige los estándares técnicos con que lo hace.

Las exigencias de la AR~4.2.2 recaen sobre los \textbf{sistemas importantes para
la seguridad}. Que un equipo lo sea depende de la \textbf{función que cumple en
la instalación}, no de cómo fue diseñado: es importante para la seguridad si su
función, o su falla, puede afectar la seguridad radiológica o nuclear. Un equipo
diseñado con el máximo rigor puede no serlo, y uno de laboratorio pasa a serlo
si se lo usa para una función de seguridad. El criterio~42 lo deja explícito:
``\emph{cuando un equipo o componente tiene varias funciones y una de las cuales
es de seguridad, debe ser clasificado como parte del sistema de
seguridad}''~\cite{ARN_AR422}.

El MCA clasifica pulsos por amplitud y entrega espectros y eventos a una PC,
Sec.~\ref{S:mca-descripcion-practica}. No actúa sobre ningún proceso, no tiene
alarmas y su salida no interviene en decisiones de operación. Por lo tanto
\textbf{no es un sistema importante para la seguridad}. Que no haya sido
diseñado para esa función es una consecuencia de esto, no la razón: tampoco está
calificado para cumplirla, como se discute en la Sec.~\ref{S:norm-iaea}.

La AR~4.2.2 ``\emph{es aplicable al diseño de reactores de investigación tipo
piscina y a las modificaciones de los sistemas importantes para la
seguridad}''~\cite{ARN_AR422}. Un instrumento de medición que no forma parte de
esos sistemas queda fuera de su alcance. Los criterios más cercanos a un MCA
marcan el límite, Tabla~\ref{t:norm-criterios-ar422}.

\begin{table}[H]
\centering
\caption[Criterios de la AR 4.2.2 más cercanos a un MCA]{%
Criterios de la AR~4.2.2 más cercanos a un MCA. A modo de comparación, el
criterio~34 exige a los sistemas de seguridad una tasa de falla por demanda
inferior a $10^{-3}$.\label{t:norm-criterios-ar422}}
\begin{tabularx}{\textwidth}{|>{\centering\arraybackslash}p{0.1\textwidth}|>{\raggedright\arraybackslash}X|>{\raggedright\arraybackslash}X|}
\hline \encabezado \textbf{Criterio} & \textbf{Qué establece} &
  \textbf{Por qué no alcanza al MCA} \\ \hline
42 & un equipo con una función de seguridad se clasifica como parte del
  sistema de seguridad & el MCA no tiene funciones de seguridad \\ \hline
50 & un sistema de calidad en el diseño que ``\emph{como mínimo, cubra los
  sistemas importantes para la seguridad}'' & el MCA no es uno de ellos \\ \hline
77 & un sistema de monitoreo de radiaciones con sensores en el reactor y
  alarmas locales & el MCA no tiene alarmas ni integra ese sistema \\ \hline
78 & tasa de falla de los sistemas de alarma y monitoreo de radiaciones
  inferior a $10^{-2}$ por demanda & se aplica al sistema del criterio~77 \\ \hline
79 a 81 & diseño de los dispositivos experimentales que se introducen en el
  reactor & tratan materiales, integridad estructural y confinamiento; el MCA
  es electrónica de lectura, afuera \\ \hline
\end{tabularx}
\end{table}

El MCA pasaría a ser importante para la seguridad si su salida alimentara una
alarma o el sistema de monitoreo de radiaciones, interviniera en decisiones de
operación, o formara parte de un sistema de protección. En ese caso aplicaría
como mínimo la meta de $10^{-2}$ por demanda del criterio~78, y el sistema de
gestión de la AR~10.6.1, obligatoria para todas las instalaciones y prácticas
reguladas~\cite{ARN_AR1061}.

Las fuentes radiactivas que se usan para probar el instrumento quedan bajo la
licencia de la institución que las tiene, según la Norma Básica de Seguridad
Radiológica AR~10.1.1~\cite{ARN_AR1011}. Esto no impone requisitos sobre el MCA.

\section{Análisis de IAEA}
\label{S:norm-iaea}

Como la ARN delega el detalle técnico, la referencia internacional es el
Organismo Internacional de Energía Atómica (IAEA). Su guía específica es la
SSG-37 (Rev.~1, 2023), sobre instrumentación, control y software importantes
para la seguridad en reactores de investigación~\cite{IAEA_SSG37}. Desarrolla
los requisitos de la SSR-3~\cite{IAEA_SSR3} y se aplica con \textbf{enfoque
graduado}: el rigor exigido es proporcional a la importancia para la
seguridad~\cite{IAEA_SSG22}. Para clasificar las funciones, la guía remite a
categorías como las de la IEC~61226~\cite{IEC61226}, que separa las funciones
en A, B y C según su importancia. Una función sin consecuencias para la
seguridad queda \textbf{no clasificada}, que es el caso del MCA. La conclusión
coincide con la de la Sec.~\ref{S:norm-arn}.

Para funciones de categoría B y C, el ciclo de vida de los circuitos programados
en HDL está definido en la IEC~62566-2~\cite{IEC62566-2}. El MCA no está
obligado a cumplirla, pero sirve como vara para ubicar lo que hicimos y lo que
no, Tabla~\ref{t:norm-ciclo-vida-62566}.

\begin{table}[H]
\centering
\caption[Ciclo de vida de la IEC 62566-2 frente a este trabajo]{%
Requisitos del ciclo de vida de la IEC~62566-2 frente a lo realizado en este
trabajo.\label{t:norm-ciclo-vida-62566}}
\begin{tabularx}{\textwidth}{|>{\raggedright\arraybackslash}p{0.25\textwidth}|>{\raggedright\arraybackslash}X|>{\raggedright\arraybackslash}p{0.22\textwidth}|}
\hline \encabezado \textbf{Requisito} & \textbf{En este trabajo} &
  \textbf{Estado} \\ \hline
especificación de requisitos & mapas de registros e interfaces de cada módulo,
  Sec.~\ref{S:ap-mca-mapa} &
  cubierto, sin trazabilidad formal requisito $\rightarrow$ prueba \\ \hline
verificación por simulación & testbenches por módulo,
  Sec.~\ref{S:intro-hw-verificacion} & cubierto \\ \hline
modelo de referencia independiente del RTL & modelo en Python comparado contra
  el RTL & cubierto \\ \hline
verificación en el dispositivo & generación interna y externa controlada,
  Sec.~\ref{S:res-etapas-validacion} & cubierto \\ \hline
validación con física real & medición con setup experimental & pendiente
  \\ \hline
verificación formal y análisis estático & --- & no realizado \\ \hline
equivalencia RTL $\leftrightarrow$ netlist & --- & no realizado \\ \hline
evaluación de componentes pre-desarrollados & placa, Linux y sistema de Red
  Pitaya usados tal como vienen, Sec.~\ref{S:hw-red-pitaya} & no evaluable
  \\ \hline
\end{tabularx}
\end{table}

La última fila deja una conclusión: la plataforma es adecuada para un
instrumento de medición, pero \textbf{no se podría calificar} para una función
de seguridad. Evaluar Linux y el sistema de Red Pitaya como software
pre-desarrollado, con el rigor de la IEC~62138~\cite{IEC62138}, no es viable.

% En la nota este bloque figura como "(sin título)", con la pregunta: "¿va en
% esta sección, con un título nuevo, o en otro capítulo?". Quedó como sección
% propia de este capítulo, con el título provisorio de la nota.
\section{Normas de medición y de equipo}
\label{S:norm-medicion}

Estas normas no son obligatorias para un prototipo de investigación, pero
definen qué tiene que medir bien un espectrómetro y cómo se informa,
Tabla~\ref{t:norm-medicion}.

\begin{table}[H]
\centering
\caption[Normas de medición y de equipo que guían el diseño del MCA]{%
Normas de medición y de equipo que guían el diseño del
MCA.\label{t:norm-medicion}}
\begin{tabularx}{\textwidth}{|>{\raggedright\arraybackslash}p{0.17\textwidth}|>{\raggedright\arraybackslash}X|>{\raggedright\arraybackslash}p{0.3\textwidth}|}
\hline \encabezado \textbf{Norma} & \textbf{Qué trata} &
  \textbf{Dónde se trata en este trabajo} \\ \hline
IEC 61452:2021~\cite{IEC61452} & calibración y uso de espectrómetros de
  germanio ($\approx$ \qty{45}{\kilo\electronvolt} a
  \qty{3}{\mega\electronvolt}); la Ed.~2 incorpora electrónica digital, Monte
  Carlo y límites de detección. Reemplaza a la ANSI/IEEE N42.14-1999,
  retirada~\cite{IEEEN4214} & Sec.~\ref{S:mca-tiempo-muerto},
  Sec.~\ref{S:mca-pileup},
  Sec.~\ref{S:implmca-perdidos},
  Sec.~\ref{S:implmca-teselado}
  \\ \hline
IEC 62755 / ANSI N42.42~\cite{IEC62755} & formato XML del espectro y sus
  metadatos (detector, calibración, tiempo real y vivo) &
  Sec.~\ref{S:sw-guardado} \\ \hline
GUM~\cite{GUM2008} & expresión de la incertidumbre de medición &
  Cap.~\ref{C:resultados} \\ \hline
IEC 61010-1~\cite{IEC61010-1} e IEC 61326-1~\cite{IEC61326-1} & seguridad
  eléctrica y compatibilidad electromagnética de equipos de laboratorio & la
  placa es un producto comercial; el trabajo cambia el bitstream y el
  software, no el hardware \\ \hline
\end{tabularx}
\end{table}

El riesgo eléctrico real del montaje es la alta tensión del detector, que no
forma parte del MCA.

% Bibliografía de este capítulo. Va al final del archivo, siempre.
\bibliografiacapitulo

%%% Local Variables:
%%% mode: latex
%%% TeX-master: "template"
%%% End:
